\chapter[Introduction to the Template: Motivation and First Steps]{Introduction to the Template Motivation and First Steps}
\label{cp:introduction}

{
\parindent0pt

\textit{Template: EPFL Thesis}

\textit{Current Version: 1.0.0}

\textit{License: \LaTeX~Project Public License v1.3c}

\textit{Official Repository: \href{https://github.com/pmdlt/epfl-thesis}{GitHub Repository}}

\vspace{.935em}

Welcome to the \textcolor{maincolor}{\textit{EPFL Thesis}} template! Thank you for choosing it for your dissertation, report, or project at the École polytechnique fédérale de Lausanne. This chapter introduces its purpose and helps you get started. See \autoref{cp:user-guide} for a detailed guide, and \autoref{cp:latex-tutorial} for a brief \LaTeX~tutorial to maximise its use.
}

\section{Motivation}
This template provides a structured, professional LaTeX foundation for EPFL students writing theses and academic reports. It combines clear file organisation, a clean design aligned with EPFL branding, and flexible class options for language, cover style, and document stage. The goal is to let you focus on your research content rather than document formatting.

\section{Getting Started}
To start using this template, you first need to know how to use \LaTeX. For this, please refer to \autoref{cp:latex-tutorial}. Once you are familiar with \LaTeX, you will need to either install it locally or use an online \LaTeX~editor.

If you prefer an online editor, \href{https://www.overleaf.com/}{Overleaf} is a popular choice. To use this template locally:

\begin{enumerate}[font=\itshape]
    \setlength{\itemsep}{.375em}
    \item \textit{Clone the \href{https://github.com/pmdlt/epfl-thesis}{repository} from GitHub.}
    \item \textit{Install a \LaTeX~distribution (TeX Live, MacTeX, or BasicTeX).}
    \item \textit{Edit \texttt{Metadata/Metadata.tex} with your thesis details.}
    \item \textit{Build with \texttt{make all} or \texttt{latexmk}.}
\end{enumerate}

If you choose to use a local editor, you must first install \LaTeX~on your machine. I recommend either \href{https://www.tug.org/texlive/}{TeX Live} or \href{https://miktex.org/}{MikTeX}. After installing \LaTeX, select an editor for writing and editing your documents. A helpful overview of editors is available on \href{https://tex.stackexchange.com/questions/339/latex-editors-ides}{TeX Stack Exchange}. Once \LaTeX~and your editor are set up, clone the template from GitHub and start writing.

\begin{block}[tip]
In this repository, you'll find a \href{https://github.com/pmdlt/epfl-thesis/blob/master/Makefile}{Makefile} and a \href{https://github.com/pmdlt/epfl-thesis/blob/master/.latexmkrc}{Latexmk configuration file}, either of which can be used to compile the project. The Makefile supports both \texttt{rubber} and \texttt{latexmk} as compilers.
\end{block}

\section{Getting Help}
As a newcomer, you may encounter situations where you want to do something in \LaTeX~or with this template but aren't sure how. When questions arise, you have several options. The \href{https://tex.stackexchange.com/}{TeX Stack Exchange} is an active community that can help with nearly any issue. The demo chapters in this document also explain how to use the template. Of course, searching online is always a reliable resource.

\subsection{Issues and Suggestions}
If you encounter a bug or have a suggestion, you can submit them via the \textit{Issues} tab in the \href{https://github.com/pmdlt/epfl-thesis}{GitHub} repository. Please be as descriptive as possible when reporting issues.

Pull requests and pushes trigger a \href{https://github.com/pmdlt/epfl-thesis/actions/workflows/latex.yml}{GitHub Action} that automatically compiles the document. If the compilation fails, the pull request checks will fail. Please keep this in mind when submitting changes.

\subsection{In-Built Comments, Guidance Texts and Warnings}
Within this document template, you may encounter informational text displaying the message ``Writing Guidance.'' These sections are provided solely as guidance to help users understand what content should be included in specific sections. They are not related to the \LaTeX~code itself within this template.

While navigating through the template, especially the configuration files, you will notice that everything is thoroughly commented. \LaTeX~can sometimes be difficult to understand without proper documentation for the packages we are working with. With that in mind, advanced modifications are accompanied by detailed comments.

Finally, regarding warnings: please take them into consideration when compiling the document. The template is designed to compile cleanly, so please strive to maintain warning-free builds.

\section{Important Notices}
This template is tailored for EPFL students writing theses, dissertations, and academic reports. Update the metadata in \texttt{Metadata/Metadata.tex} with your school, department, supervisors, and thesis details before submitting your work.

If you find this template useful, consider giving the \href{https://github.com/pmdlt/epfl-thesis}{repository} a star on GitHub to help others discover it.
